\section[几何 Satake（续）]{几何 Satake（续，2 月 1 日）}

\srcpage{128}

\subsection{回顾}

目前我们考虑分裂约化群 $G/\cO_E$。取 $T\subset B\subset G$。我们定义了 $\cY_S$ 上 $d$ 次 Cartier 除子的模空间 $\Div^d_{\cY}$，它可以看作形变了的「乘积」$S\times\Spa\cO_E$。

我们定义了 Beilinson--Drinfeld Grassmann 簇
\[
\Gr_{G,\Div^d_{\cY}}\longrightarrow\Div^d_{\cY}.
\]
$\Gr_{G,\Div^d_{\cY}}$ 参数化 $\cY_S$ 上的 $G$-主丛，连同其在 $\cY_S\setminus D_S$ 上的亚纯平凡化；或者等价地，由 Beauville--Laszlo，它参数化 $(\cY_S)^\wedge_{D_S}$ 上的对象，连同其在 $(\cY_S)^\wedge_{D_S}\setminus D_S$ 上的亚纯平凡化。特别地，$L^+G$ 通过改变平凡化作用在 $\Gr_{G,\Div^d_{\cY}}$ 上。

我们常常要考虑它上面的族，
\[
\begin{tikzcd}
\Gr_{G,S/\Div^d_{\cY}} \arrow[r] \arrow[d] & S \arrow[d] \\
\Gr_{G,\Div^d_{\cY}} \arrow[r] & \Div^d_{\cY}
\end{tikzcd}
\]
局部 Hecke 叠是 $\mathrm{Hk}_{G,S/\Div^d_{\cY}}=L^+G\backslash\Gr_{G,S/\Div^d_{\cY}}=(L^+G\backslash LG/L^+G)\times_{\Div^d_{\cY}}S$。我们定义了 $D_{\et}(\mathrm{Hk}_{G,S/\Div^d_{\cY}},\Lambda)^{bd}$。它有相对 $/S$ 的反常 $t$-结构。特别地得到阿贝尔范畴 $\Perv(\mathrm{Hk}_{G,S/\Div^d_{\cY}},\Lambda)$，它关于 $S$ 是函子性的。

\subsection{Satake 范畴上的卷积}

我们看到 $D_{\et}(\mathrm{Hk}_{G,S/\Div^d_{\cY}},\Lambda)$ 上有卷积幺半结构。我们要证明它保持相对 $/S$ 的反常 $t$-结构的核心。

\begin{proposition}
沿 $\Gr$ 的拉回诱导出全忠实函子
\[
\Perv(\mathrm{Hk}_{G,S/\Div^d_{\cY}},\Lambda)\hookrightarrow D_{\et}(\Gr_{G,S/\Div^d_{\cY}},\Lambda)^{bd}.
\]
\end{proposition}

\begin{proof}
映射 $\Gr_{G,S/\Div^d_{\cY}}\to\mathrm{Hk}_{G,S/\Div^d_{\cY}}$ 是 $L^+G$-主丛，而 $L^+G$ 是连通群。
\end{proof}

\begin{proposition}
若 $A,B\in{}^pD^{\le0}(\mathrm{Hk}_{G,S/\Div^d_{\cY}})^{bd}$，则也有 $A\star B\in{}^pD^{\le0}(\mathrm{Hk}_{G,S/\Div^d_{\cY}})^{bd}$。若 $A,B$ 是反常层，且 $A$ 是平坦反常层（即对任意 $\Lambda$-模 $M$，$A\overset{L}{\otimes}_\Lambda M$ 仍是反常层）\footnote{平坦性显然是必要的，只需看 $G$ 是平凡群的情形。}，则 $A\star B$ 仍是反常层。
\end{proposition}

\begin{proof}[证明概要]
由相对反常 $t$-结构的定义，该陈述可在几何纤维上验证。几何纤维按 untilt 分解为乘积，因此化归到 $d=1$ 以及 $A=j_{\mu_1!}\Lambda$、$B=j_{\mu_2!}\Lambda$ 的情形。

由于一切都与基变换交换，化归到 $S=\Div^1_{\cY}$。此时一切都是 ULA 的，于是可化归到特殊纤维，而该情形由 Zhu 关于 Witt 向量仿射 Grassmann 簇的工作推出 [Zhu17]。
\end{proof}

上述概略将在接下来几个小节中补全细节。

\srcpage{129}

\begin{remark}
换一种做法，我们可以使用融合积，于是不必在这里使用 Witt 向量仿射 Grassmann 簇或「卷积的半小性」。

因此我们得到平坦反常层上的卷积积。
\end{remark}

\subsection{$\mathrm{Hk}_{G,S/\Div^d_{\cY}}$ 上的 ULA 层}

\begin{definition}
称对象 $A\in D_{\et}(\mathrm{Hk}_{G,S/\Div^d_{\cY}},\Lambda)^{bd}$ 在 $S$ 上是 \emph{ULA} 的，若它拉回到 $\Gr_{G,S/\Div^d_{\cY}}$ 之后是 ULA$/S$ 的。
\end{definition}

\begin{warning}
这引入了一点不对称性。事实上，$\mathrm{Hk}_{G,S/\Div^d_{\cY}}=L^+G\backslash LG/L^+G$ 具有交换对称性，而上述定义破坏了它。结果表明「在 $S$ 上是 ULA 的」这一性质是对称的，尽管从定义中看不出这一点。
\end{warning}

\begin{proposition}
设 $d=1$。$A\in D_{\et}(\mathrm{Hk}_{G,S/\Div^d_{\cY}},\Lambda)^{bd}$ 是 ULA 的，当且仅当对一切 $\mu$ 以及 $S\xrightarrow{[\mu]}\mathrm{Hk}_{G,S/\Div^d_{\cY}}$，限制 $A|_S\in D_{\et}(S,\Lambda)$ 是局部常值的，且纤维是完美复形。
\end{proposition}

\begin{corollary}
$D_{\et}(\mathrm{Hk}_{G,S/\Div^d_{\cY}},\Lambda)^{bd}$ 中 ULA 层构成的类，在一切能由 $\mathrm{Hk}_{G,S/\Div^d_{\cY}}$ 及其分层构造出的六种运算下保持（例如 $Rj_*$、$Ri^!$、$R\mathcal{H}om$、$\ldots$）。
\end{corollary}

\begin{corollary}
考虑
\[
\begin{tikzcd}
(\Spa C)^\diamond \arrow[r] \arrow[dr] & (\Spa\cO_C)^\diamond \arrow[d] & (\Spa k)^\diamond \arrow[l] \arrow[dl] \\
& \Div^1_{\cY}=(\Spa\cO_E)^\diamond
\end{tikzcd}
\]
它诱导出（相对 $/S$ 反常 $t$-正合的）等价
\[
D^{ULA}_{\et}(\mathrm{Hk}_{G,\Spa C},\Lambda)^{bd}\xleftarrow{\ \sim\ }D^{ULA}_{\et}(\mathrm{Hk}_{G,\Spa\cO_C},\Lambda)^{bd}\xrightarrow{\ \sim\ }D^{ULA}_{\et}(\mathrm{Hk}_{G,\Spa k},\Lambda)^{bd}.
\]
特别地，反常层范畴在相应的范畴中是等价的。
\end{corollary}

\begin{proof}
在三种情形，范畴都由「严格 hensel」的分层上的层作半正交分解而构造出来，所以每个分层贡献的范畴就是 $D(\Lambda)$。这使我们可以匹配生成元，而 $\Ext$ 则由 ULA 性匹配。
\end{proof}

\begin{proof}[命题 24.6 的证明概要]
关键在于 $j_{\mu!}\Lambda$ 是 ULA 的。为验证这一点，我们过渡到 Grassmann 簇，转而考虑 $j_\mu:\Gr_{G,\mu}\hookrightarrow\Gr_{G,\le\mu}$。可以化归到泛情形，其中底为 $S=\Div^1_{\cY}$。然后只需在拉回到仿射旗簇
\[
\mathcal{F}\ell_G=LG/\operatorname{Iw}\xrightarrow{(G/B)^\diamond}\Gr_G=LG/L^+G
\]
之后验证。$\operatorname{Iw}$-作用在 $\mathcal{F}\ell_G$ 上诱导出分层
\[
\mathcal{F}\ell_G=\bigcup_{w\in\widetilde W}\mathcal{F}\ell_{G,w},
\]
其中 $\widetilde W=N(T)(B_{\mathrm{dR}})/T(B^+_{\mathrm{dR}})$ 是扩充仿射 Weyl 群。于是有短正合列
\[
1\to\underbrace{X_*(T)}_{T(B_{\mathrm{dR}})/T(B^+_{\mathrm{dR}})}\to\widetilde W\to W\to1.
\]
\end{proof}

\srcpage{130}

$\mathcal{F}\ell_{G,w}$ 的好处是它们具有很好的奇点消解。因此我们可以转而化归到 $j_{w!}\Lambda$，其中 $j_w:\mathcal{F}\ell_{G,w}\hookrightarrow\mathcal{F}\ell_G$。

于是关键在于存在 Demazure--Bott--Samelson 消解
\[
\begin{tikzcd}
& \widetilde{\mathcal{F}\ell_{G,\dot w}} \arrow[dl,"\tilde j_w"'] \arrow[dr] \\
\mathcal{F}\ell_{G,w} \arrow[rr,"j_w"'] && \mathcal{F}\ell_{G,\le w}
\end{tikzcd}
\]
其中 $w=s_{i_1}\ldots s_{i_l}\omega$，$l=\mathrm{length}(w)$，$\omega\in\Omega\subset\widetilde W$。$\Omega$ 参数化闭轨道（所以对 $\omega\in\Omega$，$\mathcal{F}\ell_{G,\omega}$ 是一个点）。我们有
\[
\widetilde{\mathcal{F}\ell_{G,\dot w}}=\mathcal{P}_{s_{i_1}}\times^{\operatorname{Iw}}\mathcal{P}_{s_{i_2}}\times^{\operatorname{Iw}}\cdots\times^{\operatorname{Iw}}\mathcal{P}_{s_{i_l}}/\operatorname{Iw},
\]
而每个 $\mathcal{P}_{s_{i_j}}/\operatorname{Iw}\cong(\mathbf{P}^1)^\diamond$，从而 $\widetilde{\mathcal{F}\ell_{G,\dot w}}$ 是迭代的 $(\mathbf{P}^1)^\diamond$-丛。这里 $\mathcal{P}_s\subset L^+G$ 是对应于简单仿射反射 $s$ 的 parahoric 子群。

这个逻辑如何运作？我们手工定义 parahoric 群，并验证所得的 Demazure--Bott--Samelson 消解是迭代的 $(\mathbf{P}^1)^\diamond$-丛。为在仿射旗簇中识别像，可以在几何点上验证，此时 $B_{\mathrm{dR}}$ 成为 $\mathbf{C}_p$ 上的 Laurent 级数环（在特征 $0$ 中）。为验证提升 $\tilde j_w$ 的存在性，也可以验证某个映射在点上同构。

由于紧合前推保持 ULA 性，只需证明 $\tilde j_{w!}\Lambda$ 在 $\widetilde{\mathcal{F}\ell_{G,\dot w}}$ 上是 ULA 的。现在 $\widetilde{\mathcal{F}\ell_{G,\dot w}}$ 是光滑的，它在 $\mathcal{F}\ell_G$ 中的边界表现得像正规交叉除子。边界由类似的胞腔组成，这些胞腔通过从 $\dot w$ 中删去因子得到。于是这给出了 $\tilde j_{w!}\Lambda$ 由分层上的常值层组成的消解，而所有分层都是光滑的。\hfill$\square$

\begin{proposition}
设 $S\to\Div^d_{\cY}$ 任意。考虑常数项函子
\[
CT_B:D_{\et}(\mathrm{Hk}_{G,S/\Div^d_{\cY}},\Lambda)^{bd}\to D_{\et}(\mathrm{Hk}_{T,S/\Div^d_{\cY}},\Lambda)^{bd}.
\]
则 $A$ 是 ULA 的，当且仅当 $CT_B(A)$ 是 ULA 的，当且仅当对投影 $\pi_T:\Gr_{T,S/\Div^d_{\cY}}\to S$，$R\pi_{T*}CT_B(A)\in D_{\et}(S,\Lambda)$ 是局部常值且取值为完美复形。
\end{proposition}

\begin{remark}
函子 $CT_B$ 可以解释为 $\mathrm{Hk}_{G,S/\Div^d_{\cY}}$ 上某个 $\mathbb{G}_m$-作用的双曲局部化，而双曲局部化一般保持 ULA 性。这里用到的另一点是 $\Gr_{T,S/\Div^d_{\cY}}\to S$ 实际上是有限的，而对有限映射，ULA 性可以在前推到基之后验证。
\end{remark}

\subsection{回到命题 24.2 的证明}

我们要证明
\[
j_{\mu_1!}\Lambda\star j_{\mu_2!}\Lambda\in{}^pD^{\le0}(\Gr_{G,\Div^1_{\cY}},\Lambda).
\]
考虑卷积 Grassmann 簇 $\widetilde{\Gr}_{G,(\Div^1_{\cY})^2}\xrightarrow{\pi_2}\Gr_{G,(\Div^1_{\cY})^2}$，它参数化 $(\cE_1,\cE_2,S^{\#}_1,S^{\#}_2)$：$\cE_1$ 在 $S^{\#}_1$ 之外的一个平凡化，以及 $\cE_1$ 与 $\cE_2$ 在 $S^{\#}_2$ 之外的一个同构。映射 $\pi_2$ 忘掉 $\cE_1$ 并复合两个同构。它在对角线之外是同构。在对角线上，它限制为卷积仿射 Grassmann 簇
\[
\widetilde{\Gr}_{G,\Div^1_{\cY}}\xrightarrow{\ \pi\ }\Gr_{G,\Div^1_{\cY}}.
\]
我们要 $R\pi_*(j_{\mu_1!}\Lambda\widetilde{\boxtimes}j_{\mu_2!}\Lambda)\in{}^pD^{\le0}(\Gr_{G,\Div^1_{\cY}})$。这整体化为 $R\pi_{2!}(j_{\mu_1!}\Lambda\widetilde{\boxtimes}j_{\mu_2!}\Lambda)$。现在，$j_{\mu_1!}\Lambda\widetilde{\boxtimes}j_{\mu_2!}\Lambda$ 在 $\widetilde{\Gr}_{G,(\Div^1_{\cY})^2}$ 上是 ULA 的。所以（紧合）前推在 $\Gr_{G,(\Div^1_{\cY})^2}$ 上是 ULA 的。我们希望它落在 ${}^pD^{\le0}$ 中。这可以在应用 $t$-正合函子 $R\pi_{T*}CT_B[\deg]$ 之后验证。

\srcpage{131}

结果落在 $D_{\et}((\Div^1_{\cY})^2,\Lambda)$ 中，并由 ULA 性而是局部常值的。在对角线之外，它由 Künneth 就是张量积，即 $R\pi_{T*}CT_B(j_{\mu_1!}\Lambda)\overset{L}{\otimes}R\pi_{T*}CT_B(j_{\mu_2!}\Lambda)\in D^{\le0}((\Div^1_{\cY})^2,\Lambda)$。由于对角线的补集是稠密的，在对角线上同样成立。

\begin{remark}
这里我们用到了「融合」的想法：把卷积重新解释为更对称的东西。
\end{remark}

\subsection{Satake 范畴}

从现在起，我们重新在 $(\Div^1_X)^d$ 上工作。注意这里我们把 $\cY$ 换成了 $X$；特别地，它生活在特征 $0$ 中。

令 $G/E$ 是任意约化群。我们将通过隐含的 étale 局部化使用先前的结果，以化归到分裂情形。

\begin{definition}
令 $I$ 是有限集，$\Lambda$ 是被某个与 $p$ 互素的 $n$ 杀死的任意环。我们定义 \emph{Satake 范畴}
\[
\mathrm{Sat}^I_G(\Lambda):=\Perv^{ULA}_{flat}\bigl(\underbrace{\mathrm{Hk}_{G,(\Div^1_X)^I}}_{\text{「}\mathrm{Hk}^I_G\text{」}},\Lambda\bigr),
\]
即 $\mathrm{Hk}^I_G$ 上平坦、ULA 的反常层 $A$ 构成的满子范畴，并带有纤维函子
\[
F:\mathrm{Sat}^I_G(\Lambda)\to\LocSys((\Div^1_X)^I,\Lambda)\xrightarrow[\mathrm{Thm}]{\ \sim\ }\Rep_{W^I_E}(\Lambda),
\]
它由 $A\mapsto\bigoplus_{i\in\mathbb{Z}}R^i\pi_{G*}A$ 给出，其中
\[
\pi_G:\underbrace{\Gr_{G,(\Div^1_X)^I}}_{\Gr^I_G}\to(\Div^1_X)^I.
\]
\end{definition}

\begin{remark}
若 $G$ 分裂，用双曲局部化可以证明 $F(A)=R\pi_{T*}CT_B(A)$。这里我们用到了：来自双曲局部化滤过（参看前文关于双曲局部化滤过的定理）的谱序列在 $F(A)$ 上由于次数集中而退化。
\end{remark}

为什么 $R\pi_{G*}$ 把 $\mathrm{Sat}^I_G(\Lambda)$ 送到 $\LocSys((\Div^1_X)^I,\Lambda)$？为什么它是正合的（且保守的）？这可以在 étale 局部验证，从而化归到 $G$ 分裂的情形。于是可以应用注 24.13 中双曲局部化的描述。我们知道常数项函子是 $t$-正合的，并且保持 ULA 性。所以 $F(A)$ 是集中在次数 $0$ 的完美复形，并且是平坦的。

\begin{remark}
把 $\LocSys((\Div^1_X)^I,\Lambda)\cong\Rep_{W^I_E}(\Lambda)$ 等同起来的定理是 Drinfeld 引理的一个版本。当 $\Lambda$ 不是有限环时，这相当微妙。
\end{remark}

\subsection{融合积}

我们定义一个函子
\[
\mathrm{Sat}^{I_1}_G(\Lambda)\times\ldots\times\mathrm{Sat}^{I_m}_G(\Lambda)\to\mathrm{Sat}^{I_1\sqcup\ldots\sqcup I_m}_G(\Lambda),
\]
使它纳入交换图
\[
\begin{tikzcd}
{*:\mathrm{Sat}^{I_1}_G(\Lambda)\times\ldots\times\mathrm{Sat}^{I_m}_G(\Lambda)} \arrow[d] \arrow[dr] \\
\mathrm{Sat}^{I_1\sqcup\ldots\sqcup I_m}_G(\Lambda) & \mathrm{Sat}^{I_1,\ldots,I_m}_G(\Lambda) \arrow[l]
\end{tikzcd}
\]
其中 $\mathrm{Sat}^{I_1,\ldots,I_m}_G(\Lambda)$ 的定义与 $\mathrm{Sat}^{I_1\sqcup\ldots\sqcup I_m}_G(\Lambda)$ 类似，但去掉 $(\Div^1_X)^I$ 中满足下述条件的轨迹：当 $i,j$ 属于不同的 $I_k$ 时有 $x_i=x_j$（这些是部分对角线）。\srcpage{132}限制 $\Gr_G^{I_1\sqcup\ldots\sqcup I_m}$ 到该轨迹，分解为 $\Gr_G^{I_1}\times\ldots\times\Gr_G^{I_m}$。所以上图的对角箭头就是外张量积。

从 $\mathrm{Sat}^{I_1\sqcup\ldots\sqcup I_m}_G(\Lambda)$ 到 $\mathrm{Sat}^{I_1,\ldots,I_m}_G(\Lambda)$ 的限制是全忠实的。这可以这样验证：应用双曲局部化，从而化归到下述陈述——$(\Div^1_X)^I$ 上局部系统的范畴全忠实地嵌入到部分对角线补集上局部系统的范畴。这很容易。

为构造上述融合积函子，只需说明对角箭头的本质像落在子范畴 $\mathrm{Sat}^{I_1\sqcup\ldots\sqcup I_m}_G(\Lambda)$ 中。为此只需造出一个具有正确限制性质的层。这可以用卷积仿射 Grassmann 簇做到。这也蕴含卷积等于融合。

把上述函子与到对角线的限制复合，我们现在得到\emph{融合积}
\[
*:\mathrm{Sat}^I_G(\Lambda)\times\ldots\times\mathrm{Sat}^I_G(\Lambda)\to\mathrm{Sat}^{I\sqcup\ldots\sqcup I}_G(\Lambda)\xrightarrow{\Delta^*}\mathrm{Sat}^I_G(\Lambda).
\]
这使每个 $\mathrm{Sat}^I_G(\Lambda)$ 成为对称幺半范畴。此外，融合积与卷积积交换。现在，只要有一个与幺半结构交换的对称幺半结构，它们必定重合（Eckmann--Hilton 论证），并与刚性幺半函子 $F$ 相容。

\begin{proposition}
函子 $F=\oplus R^i\pi_{G*}$ 是对称幺半函子\footnote{需要在交换性约束中插入一些符号。要点是 $R^i\pi_{G*}$ 应当处于次数 $i$。}
\[
\mathrm{Sat}^I_G(\Lambda)\to\Rep_{W^I_E}(\Lambda).
\]
\end{proposition}

\begin{theorem}
函子 $F:\mathrm{Sat}^I_G(\Lambda)\to\Rep_{W^I_E}(\Lambda)$ 满足 Tannaka 重构所需的一切性质，因此存在 Hopf 代数 $\mathcal{H}\in\Ind\Rep_{W^I_E}(\Lambda)$，使得
\[
\mathrm{Sat}^I_G(\Lambda)=\mathrm{CoMod}_{\mathcal{H}}(\Rep_{W^I_E}(\Lambda)).
\]
\end{theorem}

在特征 $0$ 中该陈述是清楚的。整系数情形则需要一点工作。我们需要用到 $\mathrm{Sat}^I_G(\Lambda)$ 是最高权范畴，这包括一个关于 $\Ext^2$ 群消没的非平凡陈述。

\begin{proposition}
我们有 $\mathcal{H}^I\cong\bigotimes_{i\in I}\mathcal{H}^{\{i\}}$，而 $\mathcal{H}^{\{1\}}$ 对应于一个带连续 $W_E$-作用的仿射群概形 $\check G/\Lambda$。
\end{proposition}

\begin{theorem}
存在典范同构 $\check G\cong\widehat G$；若 $\widehat G$ 的 pinning 包含分圆扭转（例如 $\Lie\widehat U_a\cong\Lambda(1)$），则它是 $W_E$-等变的。这一点下次再详细展开。
\end{theorem}

\begin{corollary}
$\mathrm{Sat}_G(\Lambda)\cong\Rep(\check G)$，并且 $\mathrm{Sat}^I_G(\Lambda)\cong\Rep(\check G^I)$，后者是在 $\Rep_{W^I_E}(\Lambda)$ 内部而言的。
\end{corollary}
